\clearpage
\section{Scalar Field}
\begin{colbox}
A scalar field \(\phi(x)\) is given, therefore under lorentz transformations it is invariant. We'll develop the equations of motion of a free real field given by the following Lagrangian density
\begin{equation}
  \lagr = \frac{1}{2}\partial^\mu\phi\partial_\mu\phi
\end{equation}
\begin{enumerate}
  \item Derive the equations of motion. What equation is this?
  \item We'll add another element to the Lagrangian
  \begin{equation}
    \lagr = \frac{1}{2}\partial^\mu\phi\partial_\mu\phi-\frac{1}{2}m^2\phi^2-\frac{\lambda}{4}\phi^4
  \end{equation}
  find the new equations of motion.
\end{enumerate}
\end{colbox}
\subsection{Equations of Motion}
Using the EL equations
\begin{align}
  \pdv{\lagr}{\phi}&=0\\
  \pdv{\lagr}{\partial_\nu\phi} &= \frac{1}{2}\ins{\pdv{\partial^\mu\phi}{\partial_\nu\phi}\partial_\mu\phi+\pdv{\partial_\mu\phi}{\partial_\nu\phi}\partial^\mu\phi}\\
  &= \frac{1}{2}\ins{\tensor{g}{^\mu^\sigma}\tensor{\delta}{_\sigma^\nu}\partial_\mu\phi + \tensor{\delta}{^\nu_\mu}\partial^\mu\phi}\\
  &= \partial^\nu\phi
\end{align}
therefore the equations of motion are
\begin{equation}
  \dalembertian \phi=0
\end{equation}
\subsection{Another Element}
The only change to the EL equations is to the element depending on the field itself
\begin{equation}
  \pdv{\lagr}{\phi} = -m^2\phi-\lambda\phi^3
\end{equation}
giving us 
\begin{equation}
  \dalembertian \phi + m^2\phi+\lambda\phi^3 =0
\end{equation}